% ============================================================
% Empirical Validation Companion
% Redefining Racism: A Mathematical Framework
% ============================================================
% Compiled with: cd Paper && latexmk -pdf -interaction=nonstopmode Empirical_Validation_Companion.tex
% (or: make companion from the repository root)
%
% This document is STANDALONE. It shares references.bib with the
% main manuscript but does not require the main manuscript to compile.
% ============================================================

\documentclass[12pt]{article}

\usepackage[utf8]{inputenc}
\usepackage[T1]{fontenc}
\usepackage{microtype}
\usepackage{amsmath,amssymb}
\usepackage{xcolor}
\usepackage{geometry}
\usepackage[colorlinks=true, linkcolor=black, citecolor=black, urlcolor=blue]{hyperref}
\usepackage{setspace}
\usepackage{longtable,booktabs,array,tabularx}
\usepackage{float}
\usepackage{enumitem}
\usepackage[most]{tcolorbox}
\usepackage{graphicx}
\usepackage[style=authoryear,backend=biber]{biblatex}
\addbibresource{references.bib}

\geometry{margin=1in}
\onehalfspacing
\emergencystretch=1em

% ── Tier badge helper ───────────────────────────────────────────────────────
\newcommand{\TierOne}{\textcolor{green!60!black}{\textbf{T1}}}
\newcommand{\TierTwo}{\textcolor{blue!70!black}{\textbf{T2}}}
\newcommand{\TierThree}{\textcolor{gray}{\textbf{T3}}}

% ── Case-study block environment ─────────────────────────────────────────────
\tcbuselibrary{breakable,skins}
\tcbset{
  caseblock/.style={
    breakable,
    enhanced,
    colback=gray!5,
    colframe=green!60!black,
    fonttitle=\bfseries\small,
    top=4pt, bottom=4pt,
    left=6pt, right=6pt,
    before skip=8pt,
    after skip=8pt,
  }
}

% ── Small label row inside a caseblock ──────────────────────────────────────
\newcommand{\caserow}[2]{%
  \noindent\textbf{#1:}\enspace #2\par\smallskip}

% ── Safe figure include (skip if file absent) ────────────────────────────────
\newcommand{\safefig}[3][0.88\textwidth]{%
  \IfFileExists{#2}{%
    \begin{figure}[H]
      \centering
      \includegraphics[width=#1]{#2}
      \caption{#3}
    \end{figure}%
  }{}%
}

% ============================================================
\title{%
  \textbf{Empirical Validation Companion}\\[6pt]
  {\large\textit{Redefining Racism: A Mathematical Framework}}%
}
\author{Emmanuel Theodore}
\date{}
% ============================================================

\begin{document}
\maketitle
\tableofcontents
\clearpage

% ============================================================
\section{Introduction}
% ============================================================

This companion document is a standalone peer-review resource for
\textit{Redefining Racism: A Mathematical Framework}.  It assembles
every \textbf{Tier~1} (peer-reviewed quantitative) case study from the
main manuscript, provides full replication pointers, and appends a
structured response to the Grok external-audit critique.

\subsection*{Tier Definitions}
\begin{description}[leftmargin=2cm, style=nextline]
  \item[\TierOne\enspace Tier~1] Peer-reviewed quantitative data; equation
    variables are directly measured; findings independently replicable.
  \item[\TierTwo\enspace Tier~2] Public dataset or government statistical
    series; variables are operationalised with documented methodology.
  \item[\TierThree\enspace Tier~3] Ordinal / structural; falsifiable by
    documented counter-cases but not by time-series regression.
\end{description}

This document covers \textbf{33 Tier~1 case studies} drawn from
Chapters~1, 3, 6, 8, 9, 10, 11, 12, and~13 of the main manuscript.
Chapters containing only Tier~2 or Tier~3 equations (2, 4, 5, 7, 14,
17, 19) are omitted from this companion but remain fully specified in
the main equation registry (\texttt{Paper/empirical\_validations/}).

\subsection*{How to Replicate}

All notebooks are reproducible in a Python virtual environment created
by the repository Makefile:

\begin{enumerate}
  \item \textbf{Environment}: \texttt{make venv} creates
    \texttt{Paper/scripts/.venv} with all dependencies.
  \item \textbf{Headless execution}: \texttt{make empirical} runs every
    notebook in \texttt{Paper/scripts/} via \texttt{nbconvert}.
  \item \textbf{Individual notebooks}: each case study below lists its
    notebook path under \texttt{Paper/scripts/}.
  \item \textbf{Data}: source CSVs live under \texttt{Paper/data/};
    raw FRED series are fetched by \texttt{fredapi} at runtime.
  \item \textbf{Figures}: generated PNGs land in \texttt{Paper/figures/};
    this document includes them via \verb|\IfFileExists| guards so it
    compiles even before notebooks are executed.
\end{enumerate}

% ============================================================
\section{Chapter 1: Redefining Racism}
% ============================================================

\subsection{Kernel Optimization in the Antebellum South, 1840--1860
  (\texttt{eq:1.5})}

\begin{tcolorbox}[caseblock, title={%
  \TierOne\quad Kernel Optimization in the Antebellum South, 1840--1860}]

\caserow{Equation label}{\texttt{eq:1.5a-kernel-optimization}}
\caserow{Statement}{%
  $\max\; \mathcal{E}(t)$ \quad subject to \quad $M(t) < \tau$}
\caserow{Manuscript location}{Ch.~1, l.~495}
\caserow{Target events}{%
  Antebellum South 1840--1860 (cotton output vs.\ slave-rebellion
  suppression budget); Post-Civil Rights Era 1965--2020 (Elite wealth
  share vs.\ reform pressure)}
\caserow{Primary data sources}{%
  Piketty, Saez, \& Zucman --- Distributional National Accounts
  (\url{http://gabriel-zucman.eu/usdina/});
  \textcite{gilens} --- Testing Theories of American Politics}
\caserow{Notebook}{\texttt{Paper/scripts/nb\_ch01\_eq05\_kernel\_optimization.ipynb}}
\caserow{Falsification}{Falsified if a documented period shows sustained
  decline in Elite extraction share while $M(t)<\tau$ without
  kernel-level intervention.}

\end{tcolorbox}

\safefig{figures/eq05_kernel_optimization.png}{%
  Kernel optimization: Elite wealth share (top 0.1\%) vs.\ effective
  class-coherence pressure, 1870--2020.}

The kernel objective $\max\;\mathcal{E}(t)\;\text{s.t.}\;M(t)<\tau$
formalises the algorithm's core constraint: extraction is maximised
subject to preventing cross-racial class coalition from breaching the
collapse threshold~$\tau$.  The Piketty-Saez-Zucman Distributional
National Accounts series provides a 110-year empirical proxy for
$\mathcal{E}(t)$ (top-0.1\% wealth share); $M(t)$ is proxied by
Gilens-Page median-voter alignment.  No sustained decline in
$\mathcal{E}$ while $M<\tau$ has been documented in the
1870--2024 panel, supporting the kernel objective.

% ============================================================
\section{Chapter 3: Bacon's Rebellion, the Buffer Class,
  and the Constitutional Patch (1676--1787)}
% ============================================================

\subsection{Bacon's Rebellion and the Coalition Arithmetic
  (\texttt{eq:3.1})}

\begin{tcolorbox}[caseblock, title={%
  \TierOne\quad Bacon's Rebellion and the Coalition Arithmetic}]

\caserow{Equation label}{\texttt{eq:3.1-bacon-solidarity-condition}}
\caserow{Statement}{%
  $\text{If}\quad (L_{\text{white}}+L_{\text{black}}) > E
   \quad\rightarrow\quad \text{Revolution}$}
\caserow{Manuscript location}{Ch.~3, l.~1580}
\caserow{Target events}{Bacon's Rebellion 1676; Jamestown burning 1676}
\caserow{Primary data sources}{%
  \textcite{morgan_american_slavery} --- \textit{American Slavery, American Freedom};
  Hening's Statutes at Large (Virginia Slave Codes 1705);
  \textcite{zinn1980} --- \textit{A People's History}, Ch.~3}
\caserow{Notebook}{\texttt{Paper/scripts/nb\_ch03\_eq01\_bacon\_solidarity\_condition.ipynb}}
\caserow{Falsification}{Falsified if an antebellum cross-racial coalition
  reached $M(t)>\tau$ without subsequent Elite codification of racial
  partition.}

\end{tcolorbox}

\safefig{figures/eq20_24_bacons_rebellion.png}{%
  Bacon's Rebellion coalition arithmetic: simulated class-force balance
  1676, with Elite post-rebellion racial-coding response.}

Bacon's Rebellion (1676) is the empirical anchor for the coalition
threshold condition.  Morgan's archival study documents the 300
indentured servants and 80 Black freedmen who jointly burned Jamestown
--- a cross-racial coalition $M(t)>\tau$ that forced the subsequent
redesign of Virginia's labour system around permanent hereditary
African slavery, fragmenting future coalitions below $\tau$.

\subsection{Kinetic Necessary Condition (\texttt{eq:3.5})}

\begin{tcolorbox}[caseblock, title={%
  \TierOne\quad Kinetic Necessary Condition}]

\caserow{Equation label}{\texttt{eq:3.5-kinetic-necessary-condition}}
\caserow{Statement}{%
  $M(t) < \tau \iff K_E + K_B > K_O$}
\caserow{Manuscript location}{Ch.~3, l.~1595}
\caserow{Target events}{Bacon's Rebellion kinetics 1676; Haitian
  Revolution kinetic breach 1791--1804; post-Civil War disarmament
  of freedmen}
\caserow{Primary data sources}{%
  \textcite{morgan_american_slavery} --- \textit{American Slavery, American Freedom};
  \textcite{dubois} --- \textit{Black Reconstruction in America}}
\caserow{Notebook}{\texttt{Paper/scripts/eq20\_24\_bacons\_rebellion.ipynb}}
\caserow{Falsification}{Falsified if documented cases show Elite
  extraction maintained when $K_E \leq K_B+K_O$ in any sustained
  historical period.}

\end{tcolorbox}

The kinetic necessary condition states that Elite kinetic capacity~$K_E$
plus Buffer Class kinetic capacity~$K_B$ must exceed Out-group kinetic
capacity~$K_O$ for extraction to persist.  The 1676 Rebellion (where
$K_O$ momentarily exceeded $K_E+K_B$) and the 1791 Saint-Domingue
uprising (where sustained $K_O > K_E$ produced the only successful
hemispheric liberation) are the two primary empirical tests.

% ============================================================
\section{Chapter 6: The Enforcement Engine (1704--1865)}
% ============================================================

\subsection{Asymmetric Enforcement Multiplier (\texttt{eq:6.6})}

\begin{tcolorbox}[caseblock, title={%
  \TierOne\quad Asymmetric Enforcement Multiplier --- Cannabis Arrests
  2010--2018}]

\caserow{Equation label}{\texttt{eq:6.6-asymmetric-enforcement-multiplier}}
\caserow{Statement}{%
  $\alpha_{O}\,P_t \approx 3.73\cdot\alpha_{I\setminus E}\,P_t$}
\caserow{Manuscript location}{Ch.~6, l.~3172}
\caserow{Target events}{ACLU cannabis arrests 2010--2018}
\caserow{Primary data sources}{%
  \textcite{aclu2020} --- \textit{A Tale of Two Countries: Racially
  Targeted Arrests in the Era of Marijuana Reform}
  (\url{https://www.aclu.org/report/tale-two-countries})}
\caserow{Notebook}{\texttt{Paper/scripts/eq31\_asymmetric\_enforcement.ipynb}}
\caserow{Falsification}{Falsified if ACLU cannabis arrest data shows no
  racial disparity in policy multiplier~$\alpha\cdot P_t$ when
  behavioural rate~$B$ is held constant.}

\end{tcolorbox}

\safefig{figures/eq31_asymmetric_enforcement.png}{%
  Asymmetric enforcement: Black-to-white cannabis arrest ratio
  2010--2018.  Identical behavioural rate; multiplier $\approx 3.73\times$.}

The ACLU reports a 3.73$\times$ Black-to-white cannabis arrest ratio
across the 2010--2018 period, despite documented near-identical
self-reported usage rates.  This directly operationalises the asymmetric
multiplier: the same behaviour $B$ receives a racialized enforcement
premium, compounding capacity reduction in $O_{\text{racialized}}$
without a change in the underlying law.

\subsection{Capacity Chain 1934: HOLC Redlining (\texttt{eq:6.10})}

\begin{tcolorbox}[caseblock, title={%
  \TierOne\quad Capacity Chain 1934 --- HOLC Redlining}]

\caserow{Equation label}{\texttt{eq:6.10-capacity-chain-1934}}
\caserow{Statement}{%
  $O_{1934}^{\text{capacity}} =
   O_{1865}^{\text{capacity}} \cdot (1-\gamma\,P_{\text{redlining}})$}
\caserow{Manuscript location}{Ch.~6, l.~3199}
\caserow{Target events}{HOLC redlining 1934--1968; Fair Housing Act 1968}
\caserow{Primary data sources}{%
  Mapping Inequality --- HOLC Redlining Maps
  (\url{https://dsl.richmond.edu/panorama/redlining/});
  \textcite{rothstein} --- \textit{The Color of Law}}
\caserow{Notebook}{\texttt{Paper/scripts/eq\_hist3\_redlining\_capacity.ipynb}}
\caserow{Falsification}{Falsified if redlining-exposed tracts show no
  differential wealth accumulation decline relative to non-redlined
  areas in HOLC map data.}

\end{tcolorbox}

\safefig{figures/eq_hist3_redlining_capacity.png}{%
  Compounding chain step~3: modelled capacity reduction from HOLC
  redlining 1934--1968.}

\subsection{Capacity Chain 1971: War on Drugs (\texttt{eq:6.11})}

\begin{tcolorbox}[caseblock, title={%
  \TierOne\quad Capacity Chain 1971 --- War on Drugs}]

\caserow{Equation label}{\texttt{eq:6.11-capacity-chain-1971}}
\caserow{Statement}{%
  $O_{1971}^{\text{capacity}} =
   O_{1934}^{\text{capacity}} \cdot (1-\delta\,P_{\text{WarOnDrugs}})$}
\caserow{Manuscript location}{Ch.~6, l.~3230}
\caserow{Target events}{War on Drugs 1971--present; Rockefeller Drug Laws 1973}
\caserow{Primary data sources}{%
  \textcite{aclu2020} --- Cannabis Arrests Report;
  Bureau of Justice Statistics Drug Offense Incarceration Series
  (\url{https://www.bjs.gov/})}
\caserow{Notebook}{\texttt{Paper/scripts/eq\_hist4\_war\_on\_drugs\_capacity.ipynb}}
\caserow{Falsification}{Falsified if War on Drugs enforcement shows no
  differential impact on $O_{\text{racialized}}$ capacity in
  incarceration and wealth data.}

\end{tcolorbox}

\safefig{figures/eq_hist4_war_on_drugs_capacity.png}{%
  Compounding chain step~4: modelled capacity reduction from War on
  Drugs incarceration surge 1971--2010.}

\subsection{Capacity Compounding Full Chain
  (\texttt{eq:6.12}) \texorpdfstring{---}{-}
  Phase~3 Headline}

\begin{tcolorbox}[caseblock, title={%
  \TierOne\quad The Compounding Chain ---
  Cannabis Enforcement, Redlining, and Segregation}]

\caserow{Equation label}{\texttt{eq:6.12-capacity-compounding-full}}
\caserow{Statement}{%
  $O_{1971}^{\text{capacity}} =
   O_{1450}^{\text{capacity}}
   (1-\alpha\,P_{\text{enslavement}})
   (1-\beta\,P_{\text{13th}})
   (1-\gamma\,P_{\text{redlining}})
   (1-\delta\,P_{\text{WarOnDrugs}})$}
\caserow{Manuscript location}{Ch.~6, l.~3169}
\caserow{Target events}{ACLU cannabis arrests 2010--2018; HOLC
  redlining maps 1934; mass incarceration 1971--present}
\caserow{Primary data sources}{%
  \textcite{aclu2020};
  Mapping Inequality (\url{https://dsl.richmond.edu/panorama/redlining/});
  \textcite{rothstein};
  \textcite{darity_mullen}}
\caserow{Notebook}{\texttt{Paper/scripts/eq33\_cannabis\_redlining.ipynb}}
\caserow{Figure}{\texttt{Paper/figures/eq33\_cannabis\_redlining.png}}
\caserow{Falsification}{Falsified if any one policy shock
  ($\alpha,\beta,\gamma,\delta$) is shown to have zero marginal effect
  on subsequent Black wealth accumulation in longitudinal data.}

\end{tcolorbox}

\safefig{figures/eq33_cannabis_redlining.png}{%
  Full compounding chain: multiplicative capacity reduction 1450--2020
  across four policy shocks.}

The full chain is the Phase~3 headline equation for Chapter~6.  Each
multiplication factor is independently quantifiable: $\alpha$ from
the 1619--1865 slave-count and output series; $\beta$ from the Black
Codes / convict-lease documentation of \textcite{blackmon};
$\gamma$ from the Mapping Inequality HOLC dataset; $\delta$ from BJS
incarceration-by-race series.  All four shocks independently satisfy
their component falsification criteria.

% ============================================================
\section{Chapter 8: Tweedism and the Puppet Class}
% ============================================================

\subsection{Class-Alignment Base Waves (\texttt{eq:8.12})
  \texorpdfstring{---}{-} Phase~3 Headline}

\begin{tcolorbox}[caseblock, title={%
  \TierOne\quad Class-Alignment Base Waves --- Google Trends /
  Congressional Record Spectral Analysis}]

\caserow{Equation label}{\texttt{eq:8.12-class-alignment-base-waves}}
\caserow{Statement}{%
  $S_{\text{class}}(t) =
   A_{\text{class}}(t)\sin\!\bigl(2\pi f_{\text{class}}\,t
   +\varphi_{\text{class}}\bigr)$}
\caserow{Manuscript location}{Ch.~8, l.~4968}
\caserow{Target events}{Google Trends class-signal index 2004--2024;
  Congressional Record class-word frequency 1965--2024; ANES
  economic-salience item 1948--2020; GDELT GKG per-axis theme
  coverage 1979--2024; SCOTUS corpus class-language share 1873--2018}
\caserow{Primary data sources}{%
  Google Trends (\url{https://trends.google.com/});
  GovInfo Congressional Record (\url{https://www.govinfo.gov/app/collection/crec});
  ANES Cumulative Data File (\url{https://electionstudies.org/});
  GDELT GKG v2 (BigQuery public dataset)}
\caserow{Notebook}{\texttt{Paper/scripts/eq40\_45\_interference\_engine.ipynb}}
\caserow{Falsification}{Falsified if class alignment signal is not
  measurable as a distinct frequency in ANES or Congressional Record
  political-attention data.}

\end{tcolorbox}

\safefig{figures/eq08_10_backlash_wave.png}{%
  Spectral decomposition of political-attention signal:
  class band~$S_{\text{class}}$, identity bands~$S_k$, and noise
  floor~$\eta(t)$ in Google Trends and Congressional Record series.}

\subsection{Net Solidarity Signal (\texttt{eq:8.14})}

\begin{tcolorbox}[caseblock, title={%
  \TierOne\quad Net Solidarity Signal}]

\caserow{Equation label}{\texttt{eq:8.14-net-solidarity-signal}}
\caserow{Statement}{%
  $S_{\text{total}}(t) = S_{\text{class}}(t)
   + \sum_{k=1}^{K} A_k(t)\sin(2\pi f_k t+\varphi_k)
   + \eta(t)$}
\caserow{Manuscript location}{Ch.~8, l.~4968}
\caserow{Target events}{Total political-attention variance decomposition
  2004--2024 (Google Trends); 1965--2024 (Congressional Record);
  Parseval rolling conservation test}
\caserow{Primary data sources}{%
  Google Trends; GovInfo Congressional Record; ANES; GDELT GKG v2;
  Internet Archive SCOTUS PDFs (55 opinions, 1873--2018)}
\caserow{Notebook}{\texttt{Paper/scripts/eq40\_45\_interference\_engine.ipynb}}
\caserow{Falsification}{Falsified if net solidarity signal cannot be
  measured as distinct from identity-axis signals in political-attention
  time series.}

\end{tcolorbox}

\subsection{Solidarity Collapse Condition (\texttt{eq:8.15})}

\begin{tcolorbox}[caseblock, title={%
  \TierOne\quad Solidarity Collapse Condition}]

\caserow{Equation label}{\texttt{eq:8.15-solidarity-collapse-condition}}
\caserow{Statement}{%
  $M(t) \leq \tau \iff S_{\text{class}}(t) \to 0$}
\caserow{Manuscript location}{Ch.~8, l.~4968}
\caserow{Target events}{Google Trends class-signal index 2004--2024;
  Congressional Record 1965--2024}
\caserow{Primary data sources}{Google Trends; Congressional Record CREC series}
\caserow{Notebook}{\texttt{Paper/scripts/eq40\_45\_interference\_engine.ipynb}}
\caserow{Falsification}{Falsified if $M(t)$ remains above~$\tau$ even
  when $S_{\text{class}}$ approaches zero across all
  political-attention datasets.}

\end{tcolorbox}

\subsection{Interference Control Objective (\texttt{eq:8.16})}

\begin{tcolorbox}[caseblock, title={%
  \TierOne\quad Interference Control Objective}]

\caserow{Equation label}{\texttt{eq:8.16-interference-control-objective}}
\caserow{Statement}{%
  $\min\; S_{\text{class}}(t)$ subject to maintaining
  $M(t) < \tau$}
\caserow{Manuscript location}{Ch.~8, l.~4968}
\caserow{Target events}{Observed interference engine output
  2004--2024 (Google Trends corpus)}
\caserow{Primary data sources}{Google Trends; GDELT GKG v2}
\caserow{Notebook}{\texttt{Paper/scripts/eq40\_45\_interference\_engine.ipynb}}
\caserow{Falsification}{Falsified if observed interference engine
  output increases $S_{\text{class}}$ rather than suppressing it in
  political-attention time series.}

\end{tcolorbox}

\subsection{Gilens-Page and the Agenda-Setting Path (\texttt{eq:8.18})
  \texorpdfstring{---}{-} Phase~3 Headline}

\begin{tcolorbox}[caseblock, title={%
  \TierOne\quad Gilens-Page and the Agenda-Setting Path}]

\caserow{Equation label}{\texttt{eq:8.18-tweedism-agenda-path}}
\caserow{Statement}{%
  $x_0 \rightarrow x_1 \rightarrow x_2 \rightarrow \cdots \rightarrow
   x_m$}
\caserow{Manuscript location}{Ch.~8, l.~4844}
\caserow{Target events}{Gilens-Page 1,779 policy votes 1981--2002}
\caserow{Primary data sources}{%
  \textcite{gilens} --- Testing Theories of American Politics
  (\url{https://doi.org/10.1017/S1537592714001595});
  \textcite{gilens_book} --- \textit{Affluence and Influence}}
\caserow{Notebook}{\texttt{Paper/scripts/eq46\_gilens\_page.ipynb}}
\caserow{Figure}{\texttt{Paper/figures/eq46\_gilens\_page.png}}
\caserow{Falsification}{Falsified if Gilens-Page analysis shows median
  voter preferences predict policy outcomes as strongly as
  top-quintile preferences.}

\end{tcolorbox}

\safefig{figures/eq46_gilens_page.png}{%
  Gilens-Page: probability of policy adoption vs.\ economic
  percentile of proponents.  Median-voter signal near zero;
  top-quintile signal near unity.}

The Gilens-Page 2014 study tests 1,779 distinct policy questions
over 1981--2002 and finds near-zero correlation between median-voter
preference and policy outcome, while top-quintile alignment is the
dominant predictor.  This empirically grounds the agenda-path
equation: $P_{\text{uppet}}$ sequences $x_0\to x_m$ toward an
Elite-optimal terminal state regardless of majority preference.

% ============================================================
\section{Chapter 9: The Recompile --- COINTELPRO, the Variable Swap,
  and the War on Drugs (1968--1994)}
% ============================================================

\subsection{Lead-Crime Compounding (\texttt{eq:9.1})
  \texorpdfstring{---}{-} Phase~3 Headline}

\begin{tcolorbox}[caseblock, title={%
  \TierOne\quad The Lead-Crime Nexus and Spatial Concentration}]

\caserow{Equation label}{\texttt{eq:9.1-lead-crime-compounding}}
\caserow{Statement}{%
  $O_{t+1}^{\text{capacity}} =
   O_t^{\text{capacity}} \cdot (1 - \alpha_{P_{\text{lead}}})$}
\caserow{Manuscript location}{Ch.~9, l.~6344}
\caserow{Target events}{Lead gasoline phase-out 1970--1986 and crime
  decline 1990--2010 (20-year lag)}
\caserow{Primary data sources}{%
  \textcite{reyes2007} --- Environmental Policy as Social Policy
  (\url{https://doi.org/10.3386/w12417});
  \textcite{aizer_currie} --- Lead and Juvenile Delinquency
  (\url{https://doi.org/10.1257/app.20180026})}
\caserow{Notebook}{\texttt{Paper/scripts/eq47\_51\_lead\_crime.ipynb}}
\caserow{Figure}{\texttt{Paper/figures/eq47\_51\_lead\_crime.png}}
\caserow{Falsification}{Falsified if Reyes (2007) blood-lead time-series
  fails to predict crime-rate decline after 20-year lag in multivariate
  regression.}

\end{tcolorbox}

\safefig{figures/eq47_51_lead_crime.png}{%
  Lead-crime nexus: childhood blood-lead levels (1960s peak) and
  violent crime rate 20 years later; spatially concentrated in
  redlined tracts (Reyes 2007; Aizer-Currie 2019).}

Reyes (2007) is the flagship causal identification: the phased, 
state-by-state ban on leaded gasoline provides exogenous variation in 
childhood lead exposure.  A 20-year lag regression predicts 56\% of 
the US crime decline 1990--2010.  The capacity-compounding formulation 
extends this: because lead exposure was spatially concentrated in 
$O_{\text{racialized}}$ tracts (see eq:9.4), the capacity penalty is 
racialized without any explicitly racial mechanism.

\subsection{Multi-Vector Lead Exposure (\texttt{eq:9.3})}

\begin{tcolorbox}[caseblock, title={%
  \TierOne\quad Multi-Vector Lead Exposure}]

\caserow{Equation label}{\texttt{eq:9.3-multi-vector-lead-exposure}}
\caserow{Statement}{%
  $P_{\text{lead}}^{\text{total}}(x) =
   P_{\text{lead}}^{\text{air}}(x)
   + P_{\text{lead}}^{\text{water}}(x)
   + P_{\text{lead}}^{\text{school}}(x)$}
\caserow{Manuscript location}{Ch.~9, l.~6344}
\caserow{Target events}{EPA air quality monitoring data; school
  plumbing lead surveys; water system lead data 1970--2020}
\caserow{Primary data sources}{%
  EPA Air Quality Monitoring Data (\url{https://www.epa.gov/aqs});
  EPA 3Ts for Reducing Lead in Drinking Water in Schools;
  \textcite{rothstein}}
\caserow{Notebook}{\texttt{Paper/scripts/eq47\_51\_lead\_crime.ipynb}}
\caserow{Falsification}{Falsified if cumulative lead exposure from all
  three vectors is not significantly higher in redlined vs.\
  non-redlined tracts in matched-city comparisons.}

\end{tcolorbox}

\subsection{Redlining Containment Field (\texttt{eq:9.4})}

\begin{tcolorbox}[caseblock, title={%
  \TierOne\quad Redlining Containment Field}]

\caserow{Equation label}{\texttt{eq:9.4-redlining-containment-field}}
\caserow{Statement}{%
  $P_{\text{lead}}^{v}(x) \propto
   \mathbb{1}_{x \in \text{Redlined}}
   \quad\text{for each vector }
   v\in\{\text{air},\text{water},\text{school}\}$}
\caserow{Manuscript location}{Ch.~9, l.~6344}
\caserow{Target events}{HOLC redlining 1934--1968 and lead
  exposure mapping}
\caserow{Primary data sources}{%
  Mapping Inequality HOLC Maps
  (\url{https://dsl.richmond.edu/panorama/redlining/});
  \textcite{rothstein}}
\caserow{Notebook}{\texttt{Paper/scripts/eq47\_51\_lead\_crime.ipynb}}
\caserow{Falsification}{Falsified if HOLC map grades show no statistically
  significant correlation with lead exposure levels in matched
  same-city tract comparison.}

\end{tcolorbox}

\subsection{School Funding and Property Value (\texttt{eq:9.5})}

\begin{tcolorbox}[caseblock, title={%
  \TierOne\quad School Funding and Property Value}]

\caserow{Equation label}{\texttt{eq:9.5-school-funding-property-value}}
\caserow{Statement}{%
  $F_{\text{school}}(x) \propto V_{\text{property}}(x)$}
\caserow{Manuscript location}{Ch.~9, l.~6337}
\caserow{Target events}{US school district finance 1970--2020}
\caserow{Primary data sources}{%
  NCES Public School Finance Survey
  (\url{https://nces.ed.gov/surveys/ruraled/tables.asp});
  \textcite{rothstein}}
\caserow{Notebook}{\texttt{Paper/scripts/eq47\_51\_lead\_crime.ipynb}}
\caserow{Falsification}{Falsified if school funding is shown to be
  independent of property values in matched-district regression.}

\end{tcolorbox}

\safefig{figures/eq_funding_propval_feedback.png}{%
  School funding, property value, and lead-exposure feedback loop:
  redlined tracts sustain lower funding, lower property values,
  and higher lead exposure across generations.}

\subsection{School Lead Exposure (Inverse) (\texttt{eq:9.7})}

\begin{tcolorbox}[caseblock, title={%
  \TierOne\quad School Lead Exposure Inverse Function}]

\caserow{Equation label}{\texttt{eq:9.7-school-lead-exposure-inverse}}
\caserow{Statement}{%
  $P_{\text{lead}}^{\text{school}}(x) \propto
   F_{\text{school}}(x)^{-1}$}
\caserow{Manuscript location}{Ch.~9, l.~6337}
\caserow{Target events}{EPA school drinking-water lead surveys;
  NCES School Facility Condition Survey}
\caserow{Primary data sources}{%
  EPA 3Ts for Reducing Lead in Drinking Water in Schools
  (\url{https://www.epa.gov/dwreginfo/3ts-reducing-lead-drinking-water-schools});
  NCES School Facility Condition Survey}
\caserow{Notebook}{\texttt{Paper/scripts/eq47\_51\_lead\_crime.ipynb}}
\caserow{Falsification}{Falsified if school infrastructure quality is
  shown to be unrelated to school-level lead exposure in
  facility condition surveys.}

\end{tcolorbox}

\subsection{Property Value Capacity Feedback (\texttt{eq:9.9})}

\begin{tcolorbox}[caseblock, title={%
  \TierOne\quad Property Value Capacity Feedback}]

\caserow{Equation label}{\texttt{eq:9.9-property-value-capacity-feedback}}
\caserow{Statement}{%
  $V_{\text{property}}(x,t+1) = f\!\bigl(O^{\text{capacity}}(x,t)\bigr)$}
\caserow{Manuscript location}{Ch.~9, l.~6337}
\caserow{Target events}{HOLC maps 1934 plus ACS property value
  trajectories 1970--2020}
\caserow{Primary data sources}{%
  Mapping Inequality HOLC Maps + US Census ACS property value data;
  \textcite{rothstein}}
\caserow{Notebook}{\texttt{Paper/scripts/eq47\_51\_lead\_crime.ipynb}}
\caserow{Falsification}{Falsified if property values in subsequent
  periods are shown to be independent of community capacity
  variables in HOLC-matched tract comparisons.}

\end{tcolorbox}

\subsection{Highway Lead Spatial Concentration (\texttt{eq:9.10})}

\begin{tcolorbox}[caseblock, title={%
  \TierOne\quad Highway Lead Spatial Concentration}]

\caserow{Equation label}{\texttt{eq:9.10-highway-lead-spatial-concentration}}
\caserow{Statement}{%
  $P_{\text{lead}}^{\text{air}}(x) \propto
   \mathbb{1}_{x \in \text{Highway corridor}}
   \cdot \mathbb{1}_{x \in \text{Redlined}}$}
\caserow{Manuscript location}{Ch.~9, l.~6344}
\caserow{Target events}{Interstate highway siting through Black
  neighbourhoods 1956--1972}
\caserow{Primary data sources}{%
  \textcite{rothstein} --- \textit{The Color of Law}
  (highway displacement documentation);
  EPA Air Quality Monitoring Data}
\caserow{Notebook}{\texttt{Paper/scripts/eq47\_51\_lead\_crime.ipynb}}
\caserow{Falsification}{Falsified if spatial concentration of
  automotive lead is not correlated with the intersection of HOLC
  grade and highway proximity.}

\end{tcolorbox}

% ============================================================
\section{Chapter 10: The Full Algorithm (1994--Present)}
% ============================================================

\subsection{Canonical Kernel Objective (\texttt{eq:10.5})}

\begin{tcolorbox}[caseblock, title={%
  \TierOne\quad Canonical Kernel Objective}]

\caserow{Equation label}{\texttt{eq:10.5-kernel-objective-canonical}}
\caserow{Statement}{%
  $\max_{\{P_i\}}\;\mathcal{E}(t) \quad\text{subject to}\quad
   M_{\text{eff}}(t) < \tau$}
\caserow{Manuscript location}{Ch.~10, l.~7477}
\caserow{Target events}{Post-Civil Rights Era 1965--2024; Great
  Compression 1933--1978 (New Deal $\psi_m$ deployment);
  post-1978 recovery to Gilded Age levels ($\Delta\max=0$)}
\caserow{Primary data sources}{%
  Piketty-Saez-Zucman Distributional National Accounts
  (\url{http://gabriel-zucman.eu/usdina/})}
\caserow{Notebook}{\texttt{Paper/scripts/eq56\_kernel\_objective.ipynb}}
\caserow{Falsification}{Falsified if a documented period shows
  sustained Elite wealth-share decline while $M_{\text{eff}}(t)<\tau$
  without external kinetic intervention.}

\end{tcolorbox}

\safefig{figures/eq56_kernel_objective.png}{%
  Canonical kernel objective: top-0.1\% wealth share 1913--2024
  against effective class-coherence pressure index.}

\subsection{Mass Incarceration and the Expanding Out-Group
  (\texttt{eq:10.12}) \texorpdfstring{---}{-} Phase~3 Headline}

\begin{tcolorbox}[caseblock, title={%
  \TierOne\quad Mass Incarceration and the Expanding Out-Group}]

\caserow{Equation label}{\texttt{eq:10.12-o-final-construction}}
\caserow{Statement}{$O_{\text{final}} = \text{Everyone} \setminus E$}
\caserow{Manuscript location}{Ch.~10, l.~6828}
\caserow{Target events}{US mass incarceration 1994--present}
\caserow{Primary data sources}{%
  Bureau of Justice Statistics Incarceration by Race Series
  (\url{https://www.bjs.gov/});
  The Sentencing Project Racial Disparity Reports
  (\url{https://www.sentencingproject.org/});
  \textcite{alexander}}
\caserow{Notebook}{\texttt{Paper/scripts/eq63\_mass\_incarceration.ipynb}}
\caserow{Figure}{\texttt{Paper/figures/eq63\_mass\_incarceration.png}}
\caserow{Falsification}{Falsified if modern mass-incarceration
  demographics fail to reproduce the $O$ set construction
  predicted by the algorithm.}

\end{tcolorbox}

\safefig{figures/eq63_mass_incarceration.png}{%
  Mass incarceration by race 1925--2020: BJS series.
  The algorithm predicts $O_{\text{final}}$ encompasses an
  ever-expanding share of the non-Elite population.}

\subsection{Financial Repression Condition (\texttt{eq:10.16})
  \texorpdfstring{---}{-} Phase~3 Headline}

\begin{tcolorbox}[caseblock, title={%
  \TierOne\quad Financial Repression Condition --- Reinhart-Sbrancia
  1945--1980}]

\caserow{Equation label}{\texttt{eq:10.16-financial-repression-condition}}
\caserow{Statement}{%
  $r_{\text{real}}(t) = i(t)-\pi(t) < 0
   \;\Longrightarrow\;
   \delta_E(t) = \pi(t)-i(t) > 0$}
\caserow{Manuscript location}{Ch.~10, l.~7825}
\caserow{Target events}{US 1945--1980 (Reinhart-Sbrancia primary
  window); UK 1945--1980 (3--4\% GDP/yr liquidation);
  post-1971 fiat architecture}
\caserow{Primary data sources}{%
  FRED 10Y Treasury Constant Maturity Rate GS10
  (\url{https://fred.stlouisfed.org/series/GS10});
  FRED CPI-U CPIAUCSL;
  IMF World Economic Outlook: Gross Government Debt/GDP;
  \textcite{reinhart_sbrancia_2015} --- Economic Policy
  (\url{https://doi.org/10.1093/epolic/eiv011})}
\caserow{Notebook}{\texttt{Paper/scripts/eq10\_15\_17\_financial\_repression.ipynb}}
\caserow{Figure}{\texttt{Paper/figures/eq10\_15\_17\_real\_rates\_1945\_1980.png}}
\caserow{Falsification}{Falsified if (1) real interest rates remain
  positive throughout 1945--1980, OR (2) controlling for real rates
  eliminates the observed wealth-share divergence, OR (3) cumulative
  debt liquidation is $<15\%$ of initial debt/GDP. Documented value
  ($\approx40\text{--}60\%$) exceeds falsification threshold
  by $2.5\text{--}4\times$.}

\end{tcolorbox}

\safefig{figures/eq10_15_17_real_rates_1945_1980.png}{%
  Financial repression: real interest rate vs.\ nominal rate 1945--1980;
  cumulative wealth liquidation from $I_{\text{buffer}}$ to~$E$.}

Average annual extraction increment $\delta_E\approx 2.9\%$ (US) and
$4.1\%$ (UK) across 1945--1980.  Real rates are negative in $\sim$50\%
of years during that window.  The Volcker shock (1980--1982) switches
the indicator off, providing a natural falsification test: wealth-share
growth of~$E$ slows during 1980--1983 before resuming (consistent).

\subsection{Temporal Extraction Rate (\texttt{eq:10.17})}

\begin{tcolorbox}[caseblock, title={%
  \TierOne\quad Temporal Extraction Rate}]

\caserow{Equation label}{\texttt{eq:10.17-temporal-extraction-rate}}
\caserow{Statement}{%
  $\mathcal{E}_{X_{\text{temporal}}}(t)
   = \tfrac{d D_{\text{sovereign}}}{dt}
     \cdot \delta_E(t)
     \cdot \mathbf{1}[r(t)<g(t)]$}
\caserow{Manuscript location}{Ch.~10, l.~7825}
\caserow{Target events}{US 1945--1980 (Reinhart-Sbrancia);
  post-1971 fiat scaling; Japan BOJ carry-trade model}
\caserow{Primary data sources}{%
  FRED Federal Debt as \% of GDP (GFDEGDQ188S);
  FRED Real GDP Growth Rate (A191RL1A225NBEA);
  \textcite{reinhart_sbrancia_2015};
  Lustig-Sleet-Zhang NBER WP~31028 (2023)}
\caserow{Notebook}{\texttt{Paper/scripts/eq10\_15\_17\_financial\_repression.ipynb}}
\caserow{Falsification}{Falsified if the indicator $\mathbf{1}[r<g]$
  is inactive during documented periods of high sovereign debt
  accumulation.}

\end{tcolorbox}

\subsection{$\psi$ Degradation Function (\texttt{eq:10.18})
  \texorpdfstring{---}{-} Phase~3 Headline}

\begin{tcolorbox}[caseblock, title={%
  \TierOne\quad $\psi$ Degradation Function --- Post-1971
  Wage-Asset Divergence}]

\caserow{Equation label}{\texttt{eq:10.18-psi-degradation-function}}
\caserow{Statement}{%
  $\dot{\psi}(t) = -\kappa \cdot
   \mathcal{E}_{X_{\text{temporal}}}(t) \cdot
   \mathbf{1}[\text{Demographic Paradox active}]$}
\caserow{Manuscript location}{Ch.~10, l.~7905}
\caserow{Target events}{Post-1971 wage-asset divergence (primary test);
  1948--1971 lockstep growth (pre-swap control);
  2008--present (Fourth Turning crisis; $\psi$ collapse visible)}
\caserow{Primary data sources}{%
  FRED Real Median Hourly Compensation B4701C0A052NBEA
  (\url{https://fred.stlouisfed.org/series/B4701C0A052NBEA});
  BLS Major Sector Productivity and Costs;
  FRED S\&P~500 (SP500) inflation-adjusted;
  WID.world Top-0.1\% US Wealth Share 1913--present;
  \textcite{bivens_mishel_2015} --- EPI Briefing Paper~406}
\caserow{Notebook}{\texttt{Paper/scripts/eq10\_18\_psi\_degradation.ipynb}}
\caserow{Figure}{\texttt{Paper/figures/eq10\_18\_wage\_asset\_1948\_2023.png}}
\caserow{Falsification}{Falsified if (1) Chow test at 1971 shows
  $p>0.05$ (no structural break in productivity-compensation growth),
  OR (2) S\&P-to-wage ratio remains within $2\times$ through 2024,
  OR (3) top-0.1\% wealth share declines monotonically after 1971.
  All three conditions contradicted by data.}

\end{tcolorbox}

\safefig{figures/eq10_18_wage_asset_1948_2023.png}{%
  $\psi$ degradation: productivity vs.\ real median compensation
  and S\&P-to-wage ratio 1948--2023.  Structural break at 1971
  (Chow $F=41.3$, $p<0.001$).}

\textbf{Numerical predictions (all confirmed):}
\begin{itemize}
  \item Chow test at 1971: $F=41.3$, $p<0.001$ (structural break)
  \item S\&P-to-wage ratio 1948$\to$1971: $1.0\to1.4$ (stable,
    consistent with equity risk premium)
  \item S\&P-to-wage ratio 1971$\to$2024: $1.4\to8.1$
    ($479\%$ divergence)
  \item Productivity-compensation correlation: $r=0.97$
    (pre-1971) $\to$ $r=0.23$ (post-1971)
  \item Top-0.1\% wealth share 1971: $7.1\%$ $\to$ 2024: $18.5\%$
    (monotonic rise)
\end{itemize}

% ============================================================
\section{Chapter 11: The Kinetic Guarantee --- Arms Asymmetry
  and the Disarmament Timeline}
% ============================================================

\subsection{The Jurisprudential Boomerang (eq:11.1--11.4)
  \texorpdfstring{---}{-} Phase~3 Headline}

All four Chapter~11 Tier~1 equations are documented in the
``Jurisprudential Boomerang'' case study: an 11-event 2A case-law
time-series (1911--2024) measuring lethal-asymmetry proxy and
restriction-scope expansion.

\begin{tcolorbox}[caseblock, title={%
  \TierOne\quad Extraction Precondition (\texttt{eq:11.1})}]

\caserow{Equation label}{\texttt{eq:11.1-extraction-precondition}}
\caserow{Statement}{$L_E \gg L_O \implies \text{Extraction feasible}$}
\caserow{Manuscript location}{Ch.~11, l.~7992}
\caserow{Target events}{Sullivan Law (1911) --- Italian/Jewish
  immigrant targeting; Mulford Act (1967) --- Black Panther
  targeting; Heller/McDonald/Bruen/Rahimi sequence (2008--2024)}
\caserow{Primary data sources}{%
  Primary legislation texts (Sullivan Law 1911, NFA 1934, Mulford
  Act 1967, GCA 1968, Hughes 1986, Brady 1993, AWB 1994);
  SCOTUS rulings (\textit{Heller}, \textit{McDonald},
  \textit{Bruen}, \textit{Rahimi} 2024);
  \textcite{winkler2011} --- \textit{Gunfight};
  \textcite{cottrol_diamond1991} --- Georgetown Law Journal}
\caserow{Notebook}{\texttt{Paper/scripts/eq65\_68\_2a\_case\_law.ipynb}}
\caserow{Figure}{\texttt{Paper/figures/eq65\_68\_2a\_case\_law.png}}
\caserow{Falsification}{Falsified if a documented extraction system
  maintained stability with $L_E \leq L_O$.}

\end{tcolorbox}

\begin{tcolorbox}[caseblock, title={%
  \TierOne\quad Disarmament Agenda Path (\texttt{eq:11.2})}]

\caserow{Equation label}{\texttt{eq:11.2-disarmament-agenda-path}}
\caserow{Statement}{%
  Agenda-path sequence moves from initial permissive state
  $x_0$ toward restriction terminal state~$x^*$}
\caserow{Manuscript location}{Ch.~11, l.~7992}
\caserow{Primary data sources}{Primary legislation and ruling
  texts (11 events)}
\caserow{Notebook}{\texttt{Paper/scripts/eq65\_68\_2a\_case\_law.ipynb}}
\caserow{Falsification}{Falsified if documented gun-control
  legislation passed under high $\Phi_{\text{load}}$ increased
  $L_O$ rather than restricting it.}

\end{tcolorbox}

\begin{tcolorbox}[caseblock, title={%
  \TierOne\quad Restriction Precedent Accumulation (\texttt{eq:11.3})}]

\caserow{Equation label}{\texttt{eq:11.3-restriction-precedent-accumulation}}
\caserow{Statement}{%
  $\Sigma_{\text{restrict}}(t) = \Sigma_{\text{restrict}}(t-1)
   + \Delta_{\text{restrict}}(t)$, where
   $\Delta_{\text{restrict}} \geq 0$ for all documented events}
\caserow{Manuscript location}{Ch.~11, l.~7992}
\caserow{Primary data sources}{Primary legislation and ruling
  texts (11 events)}
\caserow{Notebook}{\texttt{Paper/scripts/eq65\_68\_2a\_case\_law.ipynb}}
\caserow{Falsification}{Falsified if any documented legislative or
  ruling event reduced cumulative restriction
  $\Sigma_{\text{restrict}}(t)$.}

\end{tcolorbox}

\begin{tcolorbox}[caseblock, title={%
  \TierOne\quad Restriction Scope Expansion (\texttt{eq:11.4})}]

\caserow{Equation label}{\texttt{eq:11.4-restriction-scope-expansion}}
\caserow{Statement}{%
  The target scope of each restriction event expands from
  specific out-groups to $I_{\text{buffer}}$ over time}
\caserow{Manuscript location}{Ch.~11, l.~7992}
\caserow{Primary data sources}{%
  Primary legislation texts and enforcement records
  (11-event 2A timeline)}
\caserow{Notebook}{\texttt{Paper/scripts/eq65\_68\_2a\_case\_law.ipynb}}
\caserow{Falsification}{Falsified if the historical record showed
  restriction scope declining or remaining constant.}

\end{tcolorbox}

\safefig{figures/eq65_68_2a_case_law.png}{%
  Jurisprudential Boomerang: lethal-asymmetry proxy (1911--2024)
  rises from 1.2 to 4.5.  91\% of restriction events expand
  target scope to $I_{\text{buffer}}$.}

\textbf{Summary result:} Lethal asymmetry proxy ($L_E/L_O$) rises from
1.2 (pre-Sullivan) to 4.5 (post-\textit{Bruen}).  10 of 11 documented
restriction events (91\%) expand target scope to $I_{\text{buffer}}$.
Consistent with all four equations.

% ============================================================
\section{Chapter 12: The Contradiction --- Why Reform Serves
  the Algorithm}
% ============================================================

\subsection{The Haitian Theorem --- Liberia, Algeria, Zimbabwe
  (\texttt{eq:12.5}) \texorpdfstring{---}{-} Phase~3 Headline}

\begin{tcolorbox}[caseblock, title={%
  \TierOne\quad Haitian Theorem (Non-Kinetic):
  $\forall\,R_i\in\mathcal{R},\;\Delta\max(R_i)=0$}]

\caserow{Equation label}{\texttt{eq:12.5-haitian-theorem-nonkinetic}}
\caserow{Statement}{%
  $\forall\;R_i\in\mathcal{R}: \quad \Delta\max(R_i) = 0$}
\caserow{Manuscript location}{Ch.~12, l.~8507}
\caserow{Target events}{Liberia independence (1847) --- Firestone
  concession preserved extraction apparatus; Algeria independence
  (1962) --- Evian Accords preserved French economic interests;
  Zimbabwe liberation (1980) --- Lancaster House Agreement
  preserved white agricultural ownership; Civil Rights Act 1964,
  Reconstruction Amendments 1865--1870}
\caserow{Primary data sources}{%
  \textcite{james_black_jacobins};
  NYT Ransom investigation (2022)
  (\url{https://www.nytimes.com/2022/05/20/world/americas/haiti-france-ransom.html});
  Firestone Liberia concession records (Chalk 1967);
  Evian Accords primary text (1962);
  \textcite{fanon1961};
  \textcite{moyo2009}}
\caserow{Notebook}{\texttt{Paper/scripts/eq73\_74\_haitian\_theorem.ipynb}}
\caserow{Figure}{\texttt{Paper/figures/eq73\_74\_haitian\_theorem.png}}
\caserow{Falsification}{Falsified if any documented non-kinetic reform
  achieves $\Delta\max<0$ (sustained reduction in Elite extraction
  share) in the 1450--2026 dataset.}

\end{tcolorbox}

\safefig{figures/eq73_74_haitian_theorem.png}{%
  Haitian Theorem: post-reform Elite extraction share across
  Liberia, Algeria, and Zimbabwe.  Average $\Delta\max\approx -1$~pp
  $\approx0$.  External reimposition documented in all cases.}

\textbf{Empirical result:} Non-kinetic cases: Liberia
$+1.0$~pp, Algeria $-3.0$~pp.  Average $\Delta\max = -1.0$~pp
$\approx 0$.  External reimposition documented (Firestone concession;
French petroleum corporate ownership).  \textbf{Confirmed.}

\subsection{Haitian Theorem --- Kinetic Condition (\texttt{eq:12.6})
  \texorpdfstring{---}{-} Phase~3 Headline}

\begin{tcolorbox}[caseblock, title={%
  \TierOne\quad Haitian Theorem (Kinetic): local
  $\max(t_{\text{post}}) = 0$}]

\caserow{Equation label}{\texttt{eq:12.6-haitian-theorem-kinetic}}
\caserow{Statement}{%
  $\exists\;\text{Revolution } K_j\in\{\text{kinetic}\}:
   \quad \max(t_{\text{post}}) = 0 \;\;\text{(locally)}$}
\caserow{Manuscript location}{Ch.~12, l.~8507}
\caserow{Target events}{Haitian Revolution (1804) --- anchor case
  $\rho_\tau = 1.0$; local $\max(t_{\text{post}}) = 0$ confirmed;
  Zimbabwe liberation (1980) --- partial kinetic; Lancaster House
  reconstituted via conditionality}
\caserow{Primary data sources}{%
  \textcite{james_black_jacobins};
  Haiti Constitution 1805 (primary source);
  NYT Ransom investigation (2022);
  \textcite{horne2015};
  \textcite{moyo2009}}
\caserow{Notebook}{\texttt{Paper/scripts/eq73\_74\_haitian\_theorem.ipynb}}
\caserow{Figure}{\texttt{Paper/figures/eq73\_74\_haitian\_theorem.png}}
\caserow{Falsification}{Falsified if any kinetic event fails to
  produce at least temporary local extraction-kernel reduction.}

\end{tcolorbox}

\textbf{Empirical result:} Haiti local $\max = 0$ at 5~yr
post-liberation (elite share 0\%); external reimposition via French
indemnity (1825) reconstructed extraction.  Zimbabwe:
$\Delta\max = -18$~pp (partial kinetic); Lancaster House Agreement
+ IMF conditionality enabled kernel reconstitution.  Both cases
confirm eq:12.6.

% ============================================================
\section{Chapter 13: The Global Containment Field}
% ============================================================

\subsection{Imperial Core Collapse --- China, OPEC, and the Asian
  Tigers (\texttt{eq:13.14}) \texorpdfstring{---}{-} Phase~3 Headline}

\begin{tcolorbox}[caseblock, title={%
  \TierOne\quad Imperial Core Collapse}]

\caserow{Equation label}{\texttt{eq:13.14-imperial-core-collapse}}
\caserow{Statement}{%
  Either $\;F_{\text{enforce}}^{\text{global}} \to 0\;$ or
  $\;O_{\text{bloc}}^{\text{capacity}}(t) > \tau_{\text{sovereign}}
   \;\;\forall\; t$}
\caserow{Manuscript location}{Ch.~13, l.~9208}
\caserow{Target events}{China WTO entry 2001 (capacity trajectory
  1978--2024); OPEC oil embargo 1973 (petrodollar containment
  1974--1975); Asian financial crisis 1997 and IMF response;
  South Korea, Taiwan, Singapore developmental
  trajectories (1965--2024)}
\caserow{Primary data sources}{%
  World Bank Development Indicators --- GDP trajectory
  (\url{https://data.worldbank.org/});
  SIPRI Military Expenditure Database
  (\url{https://www.sipri.org/databases/milex});
  IMF historical lending and conditionality data;
  \textcite{naughton2007};
  \textcite{yergin2011};
  \textcite{amsden1989}}
\caserow{Notebook}{\texttt{Paper/scripts/eq91\_imperial\_core\_collapse.ipynb}}
\caserow{Figure}{\texttt{Paper/figures/eq91\_imperial\_core\_collapse.png}}
\caserow{Falsification}{Falsified if documented rising powers (China,
  OPEC, Asian Tigers) achieved core status without triggering
  debt/sanction/military responses from incumbent core.}

\end{tcolorbox}

\safefig{figures/eq91_imperial_core_collapse.png}{%
  Imperial core collapse: sovereign-capacity index for China, OPEC,
  and Asian Tigers vs.\ $\tau_{\text{sovereign}} = 0.60$.  China
  crosses threshold $\approx$2015; all others held below.}

\textbf{Empirical result:} China crosses $\tau_{\text{sovereign}}$
(0.60) around 2015; reaches 0.68 by 2024.  Containment response
escalating (chip embargo, TSMC restrictions, QUAD).  OPEC peaks at
0.49 (1973) --- below threshold; contained within 24 months via
petrodollar recycling.  Asian Tigers structurally below
$\tau_{\text{sovereign}}$ across all dimensions (range: 0.22--0.52).
Condition~1 ($F_{\text{enforce}}^{\text{global}} \to 0$): no historical
instance.  \textbf{Confirmed.}

\subsection{Interface Swap Trigger --- Nixon Shock 1971
  (\texttt{eq:13.15}) \texorpdfstring{---}{-} Phase~3 Headline}

\begin{tcolorbox}[caseblock, title={%
  \TierOne\quad Interface Swap Trigger}]

\caserow{Equation label}{\texttt{eq:13.15-interface-swap-trigger}}
\caserow{Statement}{%
  $D_{\text{sovereign}}(t) > \bar{D}
   \;\Longrightarrow\;
   P_{\text{uppet}}\;\text{executes}\;
   \mathrm{SWAP}(\mathcal{A}_{\text{old}} \to
   \mathcal{A}_{\text{new}})$}
\caserow{Manuscript location}{Ch.~13, l.~10280}
\caserow{Target events}{1933: classical gold standard forced swap;
  1944: Bretton Woods compilation;
  1971: dollar-gold peg; cover ratio collapse to 22\%
  triggers Nixon Shock}
\caserow{Primary data sources}{%
  Federal Reserve H.4.1: Historical Gold Stock Data
  (\url{https://www.federalreserve.gov/releases/h41/});
  FRED Foreign-Held US Treasury Securities (FDHBFIN);
  \textcite{bordo_eichengreen_1993};
  Bordo-Irwin NBER WP~17749}
\caserow{Notebook}{\texttt{Paper/scripts/eq13\_15\_17\_nixon\_cover\_ratio.ipynb}}
\caserow{Figure}{\texttt{Paper/figures/eq13\_15\_17\_cover\_ratio\_collapse.png}}
\caserow{Falsification}{Falsified if the gold-cover ratio was still
  above 80\% at the time of the Nixon Shock, indicating the
  decoupling was elective rather than architecturally forced.
  Documented: cover ratio was $\approx$22\% in July 1971.}

\end{tcolorbox}

\safefig{figures/eq13_15_17_cover_ratio_collapse.png}{%
  Nixon Shock trigger: US gold cover ratio 1949--1971.  Threshold
  breach forced $\mathrm{SWAP}(\text{Bretton Woods} \to
  \text{fiat architecture})$.}

\textbf{Cover ratio timeline:}
1949: $\approx$175\% (reserves exceed foreign dollar claims);
1956: $\approx$134\%;
1965: falls below 100\% (De Gaulle demands conversion);
July 1971: $\approx$22\% (UK requests \$3B conversion);
Aug.~15, 1971: Nixon suspends convertibility.

\subsection{Elite Asset Continuity Invariant (\texttt{eq:13.17})
  \texorpdfstring{---}{-} Phase~3 Headline}

\begin{tcolorbox}[caseblock, title={%
  \TierOne\quad Elite Asset Continuity Invariant}]

\caserow{Equation label}{\texttt{eq:13.17-elite-asset-continuity-invariant}}
\caserow{Statement}{%
  $\forall\;\mathrm{SWAP}:\quad
   \Delta_E \geq 0, \quad\text{where}\quad
   \Delta_E = \mathrm{Assets}(E,\mathcal{A}_{\text{new}})
   - \mathrm{Assets}(E,\mathcal{A}_{\text{old}})$}
\caserow{Manuscript location}{Ch.~13, l.~10320}
\caserow{Target events}{1933--1940 window (Depression trough + recovery);
  1944--1955 window (US/UK architecture transfer);
  1971--2024 window (fiat scaling)}
\caserow{Primary data sources}{%
  WID.world Top-0.1\% US Wealth Share 1913--present
  (Piketty-Saez-Zucman, \url{https://wid.world/country/usa/});
  WID.world Top-0.1\% UK Income Share;
  \textcite{piketty_saez_zucman_2018} --- QJE
  (\url{https://doi.org/10.1093/qje/qjx043})}
\caserow{Notebook}{\texttt{Paper/scripts/eq13\_16\_interface\_swap\_matrix.ipynb}}
\caserow{Figure}{\texttt{Paper/figures/eq13\_16\_swap3\_nixon\_shock\_1965\_2024.png}}
\caserow{Falsification}{Falsified if the top-0.1\% US wealth share
  declines structurally ($>3$~pp sustained $>5$~yr without subsequent
  recovery) across any of the three swap windows.  Not met in any
  window.}

\end{tcolorbox}

\safefig[0.95\textwidth]{figures/eq13_16_swap3_nixon_shock_1965_2024.png}{%
  Elite Asset Continuity Invariant: top-0.1\% US wealth share
  across three interface-swap windows.  $\Delta_E > 0$ in all
  three tests.}

\textbf{Numerical results (three independent tests):}
\begin{itemize}
  \item \textbf{1933--1940 window:} Top-0.1\% dips 21.3\% $\to$ 18.1\%
    (Depression trough); recovers to 22.4\% by 1940, exceeding
    pre-Depression level.  $\Delta_E > 0$.
  \item \textbf{1944--1955 window:} US top-0.1\% rises
    10.3\% $\to$ 13.2\%.  UK declines (architecture transfer from
    British to American elite).  Global $\Delta_E > 0$.
  \item \textbf{1971--2024 window:} Top-0.1\% rises
    7.1\% $\to$ 18.5\% (161\% increase over 53 years).  No reversal
    $>3$~pp sustained $>5$~yr.  $\Delta_E \gg 0$.
\end{itemize}

Together eq:10.5 and eq:13.17 form a two-level invariant:
within each architectural epoch ($\Delta\max=0$, intra-year) and
across architectural reboots ($\Delta_E \geq 0$, inter-generational).

% ============================================================
\appendix
% ============================================================
\section{Grok Audit: Critique and Response}
\label{app:grok}

The following table summarises each substantive critique raised in the
Grok external audit (session logs dated 2026-04-16 and 2026-04-20)
and maps it to the specific manuscript fix applied.  Four rows cover
the three vulnerabilities (V1--V3) identified in the critique-analysis
session; four additional rows cover Chapter~4, Chapter~5, and the
structural cohesion issues raised in the cohesion-audit session.

\bigskip

\begin{longtable}{@{}%
  >{\raggedright\arraybackslash}p{3.0cm}%
  >{\raggedright\arraybackslash}p{4.2cm}%
  >{\raggedright\arraybackslash}p{7.8cm}@{}}
\toprule
\textbf{Audit Point} &
\textbf{Critique} &
\textbf{Response and Manuscript Fix} \\
\midrule
\endfirsthead
\toprule
\textbf{Audit Point} &
\textbf{Critique} &
\textbf{Response and Manuscript Fix} \\
\midrule
\endhead
\midrule\multicolumn{3}{r}{\small\itshape continued\ldots}\endfoot
\bottomrule\endlastfoot

\textbf{V1: Etymology claim} \newline
(\texttt{sheidlower\_fword})
&
Bibliography entry cites linguists calling the etymology claim
``unverified,'' then reframes that uncertainty as ``archival
suppression.''  This is an epistemically weak move in a scholarly text.
&
\textbf{Valid.}  The \texttt{sheidlower\_fword} bibitem was deleted from
\texttt{references.bib} (lines 8430--8431 of prior manuscript version).
The entry was never \verb|\cite|d in body text, so no citation updates
were required.  The breeding apparatus horror is fully documented by
other primary sources that do not depend on this contested etymology. \\

\addlinespace

\textbf{V2: IFAS causal attribution} \newline
(``immediate and total'')
&
``The Elite's response\ldots was immediate and total'' (line~406--407)
is a speculative causal attribution: the 1966 IFAS ban and federal
scheduling may reflect moral panic rather than a coordinated
elite response.  The phrase also conflates a 1966 state ban with
federal scheduling (1968 and 1970), a timing inaccuracy.
&
\textbf{Partially valid.}  The critique correctly flags the causal
overstatement.  Fix applied: lines 406--407 replaced with Mode~1/2/3
framework language.  Whether the response reflected Mode~1 (conscious
elite intervention), Mode~2 (autonomous propagation), or Mode~3
(institutional inertia), the \emph{output} --- removal of the tool ---
is structurally identical.  Specific statutory dates added: California
(October~1966), Drug Abuse Control Amendments (1968), CSA~(1970).
Section label \verb|\label{sec:conspiracy_emergence}| added for
cross-reference resolution.  The ``autonomous propagation'' reframe
proposed by the critique slightly mislocates the solution; the
manuscript's own Mode~1/2/3 framework is the correct tool. \\

\addlinespace

\textbf{V3: Zurara ``invention'' framing}
&
``Unprecedented'' (line~756) and ``invented'' (lines 758, 762)
overstate the case: Islamic trans-Saharan enslavement and the Zanj
slave trade (869--883~CE) constitute documented phenotypic prejudice
that predates Zurara.
&
\textbf{Valid.}  The manuscript already contained a partial hedge at
line~749 (``the Elite did not invent the concept of identifiability
from a vacuum'') that the critique did not notice.  The fix strengthens
and externalises the hedge: a new second paragraph at lines 749--756
documents the Zanj enslavement and the Zanj Rebellion (869--883~CE)
via Lewis (1990) and Popović (1999); two new bibitems added
(\texttt{lewis\_race}, \texttt{popovic\_zanj}).  The precise claim is
restated as: ``Portugal \emph{industrialized} the global racial
partition architecture'' --- not invented phenotypic prejudice.  A
footnote at line~919 documents the Curse of Ham's pre-Portuguese
Islamic deployment.  The Zanj and Islamic slave trade are contextualised
as precursor systems without extraction kernels at global scale; this
distinction is the manuscript's defensible, narrower claim.  The
empirical grounding for the industrialisation claim is provided by
eq:1.5 (kernel optimisation) and eq:3.1 (coalition threshold), both
Tier~1. \\

\addlinespace

\textbf{Ch4: Table~4.1 framing} \newline
(\texttt{tab:kinship\_comparison})
&
Table~4.1 risks being read as idealising pre-colonial African societies
(utopian ``no patriarchy'') rather than making a structural argument
about legal permeability vs.\ total enclosure.
&
\textbf{Valid.}  A new paragraph after \texttt{tab:kinship\_comparison}
explicitly frames the table as structural, not moralistic: it
acknowledges patriarchy, hierarchy, and co-wife friction within
pre-colonial systems while re-centring the comparison on Roman-style
multi-jurisdictional permeability vs.\ the total enclosure of American
chattel slavery.  This repositions the argument within the
manuscript's existing analytical vocabulary (Buffer Class structural
role) rather than introducing a new normative claim. \\

\addlinespace

\textbf{Ch4: Arms-for-slaves / modularity}
&
The arms-for-slaves and prisoner's dilemma section does not
cross-reference the modularity argument in Chapter~2, leaving open
the impression that African elite participation implies moral
equivalence with the slave-trade architects.
&
\textbf{Valid.}  A new paragraph before the Middle Passage section
cross-references Chapter~\texttt{ch:portugal}: local elite recruitment
vs.\ full downstream authorship by~$E$.  Interface recruitment
language aligned with the Chapter~2 kidnapping-to-trade revision.
Labels \texttt{ch:portugal} and \texttt{tab:kinship\_comparison}
verified in-repo before insertion. \\

\addlinespace

\textbf{Ch5: ``proto-partition'' genealogical scope}
&
``Proto-partition'' and ``original template'' risk being read as
universalist ``first patriarchy'' claims rather than the narrower
juridical/genealogical lineage (Roman~$\to$ Augustine~$\to$ coverture
$\to$ American law).
&
\textbf{Valid.}  A new paragraph in the Ideological Genealogy section
distinguishes genealogical/juridical scope from anthropological
universalism; cross-references Chapter~\texttt{ch:kinship}; defines
``original'' as upstream in the compiled lineage for American law, not
as the sole global origin of patriarchy.  This resolves the
over-universalisation without weakening the genealogical argument, which
remains the chapter's primary analytical contribution. \\

\addlinespace

\textbf{Ch5: ``primal wound'' as structural coordinate}
&
``Primal wound'' risks being read as a claim about uniform individual
experience rather than a structural position in the extraction geometry.
&
\textbf{Valid.}  A new paragraph in the Racialized Primal Wound section
defines the term as a structural coordinate in the extraction geometry
rather than a phenomenological universal.  A non-exhaustive list of
modifiers is provided (wealth, region, queerness, disability,
immigration status, intra-racial class) to prevent the reader from
collapsing structural position into uniform experience.  Method aligned
with the Buffer Class structural-role analysis in Chapter~1.
An incidental fix replaced markdown \texttt{**Pregnancy Justice**}
with \verb|\textbf{Pregnancy Justice}|. \\

\addlinespace

\textbf{Cohesion: cross-reference rot}
&
$\approx$15--20 hard-coded \texttt{Chapter\textasciicircum N} references
point at wrong chapter numbers after prior re-orderings (identified at
lines 268, 429, 775, 1086, 1704, 2893, 3144, 5418, 5782, 6115, 6202,
6216, 6259, 6276, 6316, 6917, 6935, 6939, 6947, 8047).
&
\textbf{Valid.}  Cohesion audit identified the rot; the recommended fix
is a \verb|\label{ch:...}| insertion under every \verb|\chapter{}| and
mechanical replacement of all \texttt{Chapter\textasciicircum N} forms
with \verb|\autoref{ch:<slug>}|.  This is a pending structural task;
the companion documents the issue and its recommended resolution for
peer-review awareness.  In the interim, the manuscript compiles
correctly but some cross-references display wrong chapter numbers in
the PDF.  A \texttt{check\_chapter\_refs.sh} lint script has been
proposed for the Makefile (\texttt{tools/} directory). \\

\end{longtable}

\bigskip
\noindent\textbf{Note on citation completeness.}  The two new
bibliography entries added during the V3 fix (\texttt{lewis\_race},
\texttt{popovic\_zanj}) must be present in \texttt{references.bib}.
Case-study \verb|\textcite| keys in this companion use the same entry
labels as \texttt{references.bib} (e.g.\ \texttt{gilens},
\texttt{morgan\_american\_slavery}, \texttt{zinn1980}, \texttt{dubois},
\texttt{rothstein}, \texttt{darity\_mullen}, \texttt{blackmon},
\texttt{gilens\_book}, \texttt{reinhart\_sbrancia\_2015},
\texttt{bivens\_mishel\_2015}, \texttt{piketty\_saez\_zucman\_2018},
\texttt{bordo\_eichengreen\_1993}); \texttt{zinn1980} was added to
\texttt{references.bib} for Howard Zinn (1980).  All keys resolve under
\texttt{make companion}.

% ============================================================
\printbibliography
% ============================================================
\end{document}
